\documentclass[10pt,a4paper]{article}
\usepackage{amsmath,amssymb,amsthm,graphicx}
\usepackage{hyperref}
\usepackage{enumitem}
\usepackage{mathtools}
\DeclareMathOperator{\supp}{supp}
\DeclareMathOperator{\vars}{vars}
\DeclareMathOperator{\Qdeg}{Qdeg}
\DeclareMathOperator{\sing}{sing}
\DeclareMathOperator{\nsing}{nsing}
\DeclareMathOperator{\Split}{Split}
\title{Exponential polynomial-calculus size lower bounds persist under many local extension variables over finite fields}
\author{Russell Impagliazzo, Sasank Mouli, Toniann Pitassi}
\date{doi:10.4230/LIPIcs.CCC.2023.7; CC BY 4.0}
\begin{document}
\maketitle
\noindent\textit{Editorial note.} Editable author proof transcribed from official CC BY4 published PDF; all necessary author statements, local support proofs, Algorithm1/2 and Appendix retained. Original literal lemma/algorithm references and author expressions are preserved, including author-internal numbering inconsistencies. No theorem repair, new assumption, proof certification or extra counted support result is introduced. Display formulas use line breaks for page width. All production and validation run through SSH on Yupeng.\par\medskip
\noindent\textit{Allowed standard prerequisites (editorial).}\begin{itemize}
\item Polynomial-calculus derivational completeness and standard size-degree tradeoffs over Boolean bases. Exact original authors use retained, including p=2 case in Section3; standard results in references6/10/17 not recursively reproved
\item Degree hardness of boundary-expanding parity CSPs. Original AppendixLemma35 explicitly invokes reference1,Theorems3.8/4.4. Source-local random expansion and unsatisfiability constructions are fully retained.
\end{itemize}\medskip
\paragraph{Original parameter regimes (introduction, Theorems1/2).}
\textbf{Theorem1 (high-end).} For $n$ sufficiently large, there is a family of CNF tautologies $F^{\mathsf{SEL}}$ on $O(n\log n)$ variables with $\operatorname{poly}(n)$ clauses of width $O(\log n)$ such that for any $M=n^{\operatorname{polylog}(n)}$ and $\kappa=O(\log\log n)$, any PC refutation over $\mathbb F_p$ of $F^{\mathsf{SEL}}$, together with $M$ $\kappa$-local extension axioms, requires size $2^{\Omega(n/\operatorname{polylog}(n))}$.
\textbf{Theorem2 (low-end).} For the same family of tautologies as above, there are $0<\alpha,\beta,\gamma<1$ so that, for $M=n^{1+\alpha}$, $\kappa=\beta\log n$, any PC refutation of $F^{\mathsf{SEL}}$ together with any $M$ $\kappa$-local extensions over $\mathbb F_p$ requires size $2^{\Omega(n^{\gamma})}$.
We remark that our extension variables are only allowed to depend on the original variables, and not on previously defined extension variables. (In the more general case where extension variables are defined recursively, the proof system corresponds to $\mathsf{AC}^0[p]$-Frege, where the level of recursion corresponds to the $\mathsf{AC}^0[p]$ circuit depth.) Thus our lower bound can be (roughly) seen as proving exponential lower bounds for the following restricted class of depth-2.5 PC refutations. First, the refutation is given a new set of $M$ variables, $z_1,\ldots,z_M$, and is allowed to define a corresponding set of $M$ $\kappa$-local polynomials $Q_1,\ldots,Q_M$ (where each $Q_i$ can only depend on $\kappa$ original variables). Lines in the refutation are polynomials over the original variables, plus the new extension variables (which are placeholders for the $Q_i$'s). Substituting the $Q_i$'s for the new variables gives a set of depth2.5 algebraic circuits using a pre-specified set of $\kappa$-local functions at the bottom layer of the circuit.

\section*{2 Preliminaries}
\paragraph{Definition3 (Polynomial Calculus/Polynomial Calculus Resolution).}
Let $\Gamma=\{P_1\ldots P_m\}$ be an unsolvable system of polynomials in variables $\{x_1\ldots x_n\}$ over $\mathbb F$. A PC (Polynomial Calculus) refutation of $\Gamma$ is a sequence of polynomials $\{R_1\ldots R_s\}$ such that $R_s=1$ and for every $\ell\in[s]$, $R_{\ell}\in\Gamma$, $R_{\ell}$ is either a polynomial from $\Gamma$, or is obtained from two previous polynomials $R_j,R_k$, $j,k<\ell$ by one of the following derivation rules:
\[R_{\ell}=\alpha R_j+\beta R_k\text{ for }\alpha,\beta\in\mathbb F,\qquad R_{\ell}=x_iR_k\text{ for some }i\in[n].\]
The size of the refutation is $\sum_{\ell=1}^s|R_{\ell}|$, where $|R_{\ell}|$ is the number of monomials in the polynomial $R_{\ell}$. The degree of the refutation is $\max_{\ell}\deg(R_{\ell})$.
A PCR (Polynomial Calculus Resolution) refutation is a PC refutation over the set of Boolean variables $\{x_1\ldots x_n,\bar x_1\ldots\bar x_n\}$ where $\{\bar x_1\ldots\bar x_n\}$ are twin variables of $\{x_1\ldots x_n\}$ i.e. the equations $x_i^2-x_i=0$, $\bar x_i^2-\bar x_i=0$ and $x_i\bar x_i=0$ are treated as axioms.
\paragraph{Definition4 (PC plus Extension Axioms).}
Let $\Gamma=\{P_1\ldots P_m\}$ be a set of polynomials in variables $\{x_1\ldots x_n\}$ over a field $\mathbb F$. We will refer to the polynomials in $\Gamma$ as (initial) axioms. Let $\mathbf z=z_1\ldots z_M$ be new extension variables with corresponding extension axioms $z_j-Q_j(x_1\ldots x_n)$. A PC+Ext (PC plus extension) refutation of $\Gamma$ with $M$ extension axioms $\mathsf{Ext}=\{z_j-Q_j(x_1\ldots x_n)\}$ is a PC refutation of the set of polynomials $\Gamma'=\{P_1\ldots P_m,z_1-Q_1\ldots z_M-Q_M\}$. An extension axiom $z_j=Q_j(x_1\ldots x_n)$ is $\kappa$-local if $Q_j$ is a $\kappa$-junta; that is, if the polynomial $Q_j$ defining $z_j$ involves at most $\kappa$ of the $x$-variables. We say that $\Pi$ is a $(M,\kappa)$-PC+Ext refutation of $\Gamma$ if it is a PC+Ext refutation of $\Gamma$ with $M$ extension axioms, each of which are $\kappa$-local. The size of the refutation is total size of all lines in the refutation, including the polynomials in $\Gamma$ plus the extension axioms (where the size of a line $P\in\Pi$ is the number of monomials in $P$).
We note that our definition of extension axioms is more limited than the general notion of extension axioms. Here we only allow the extension variables to depend on the original variables from $\Gamma$; the more general definition allows the extension variables to depend on the original $x$-variables, and also on other extension variables.
\paragraph{Definition5 ($k$-local CSPs).}
A constraint $C_i$ over Boolean variables $\{x_1,\ldots,x_n\}$ is simply a Boolean formula over these variables. $C_i$ is a $k$-local constraint if $C_i$ depends on at most $k$ variables. A $k$-CSP $C=C_1\land\ldots\land C_m$ over $\{x_1,\ldots,x_N\}$ is the conjunction of a set of $k$-local constraints.
We translate a $k$-CSP formula into a system of polynomial equations using the standard PCR translation which we define next.
\paragraph{Definition6 (Converting $k$-CSPs into Polynomial Equations).}
Let $C$ be a $k$-local constraint over variables $x_{i_1},\ldots,x_{i_k}$. We convert $C$ to a polynomial equation, $p(C)$, using the translations $p(x_{i_j})=1-x_{i_j}$, $p(\neg A)=1-p(A)$, $p(A\lor B)=p(A)\cdot p(B)$. It is easy to check that for any Boolean assignment $\alpha$ to the underlying variables, $C(\alpha)=1\leftrightarrow p(\alpha)=0$, and $C(\alpha)=0\leftrightarrow p(\alpha)=1$.
A $k$-CSP $C=C_1\land\ldots\land C_m$ over $\{x_1,\ldots,x_n\}$ converts to a set of polynomial equations $\{E_j\mid j\in[m]\}\cup\{B_i\mid i\in[n]\}$ over $\{x_1,\ldots,x_n\}\cup\{\bar x_1,\ldots,\bar x_n\}$ where $E_j$ is the polynomial equation $p(C_j)$. In addition, we add the Boolean axioms $\{B_i\mid i\in[n]\}$, where $B_i=\{x_i^2-x_i=0,\bar x_i^2-\bar x_i=0,x_i\bar x_i=0\}$ which force $x_i,\bar x_i$ to be zero-one valued, and force exactly one of $x_i,\bar x_i$ to be one.
\section*{3 The Hard Formulas}
We distinguish between the case $p=2$ and the case $p>2$, and concentrate on the latter. This is because the case $p=2$ does not require any new technical ideas, and we can pick from a large number of known hard tautologies for this case, such as random CNF's. Over $\mathbb F_2$, every extension variable is zero-one valued, and so standard size-degree tradeoffs pertain even with respect to extension variables. Also, $\kappa$-local extension variables can change the degree by at most a factor of $\kappa$, therefore a degree lower bound of $\Omega(n)$ for the original tautology over $n$ variables implies a degree lower bound of $\Omega(n/\kappa)$ after adding $\kappa$-local extension variables. Known size-degree tradeoffs imply that the degree must be at least square root of the number of variables in order to obtain exponential size lower bounds, this immediately gives a lower bound tolerating close to $n^2/\kappa^2$ many $\kappa$-local extension variables [10,6,17].
Over any field, there are unsatisfiable families of $k$-CNF formulas (e.g. the Tseitin tautologies as well as random parity equations) that require linear degree but have polynomial sized proofs with a linear number of extension variables [8,6]. Therefore formulas that require high PC degree are not sufficient. Instead we will create our hard examples by taking a hard instance and then using selector variables to pick out a subset of the constraints. Similar ideas were used earlier (e.g., [12]). In more detail, our underlying hard unsatisfiable formulas, $\{F^{\mathsf{SEL}}_{n,k}\}$, will be constructed from a family of $k$-CSP formulas, $F_{n,k}$, that have the property that any sufficiently large subset of the constraints of $F_{n,k}$ is unsatisfiable and still requires large PC degree.
\paragraph{Definition7.}
Let $F_{n,k}=\{E_j\mid j\in[m]\}\cup\{B_i\mid i\in[n]\}$ be the system of degree-$k$ polynomial equations over $\mathbf x=\{x_i,\bar x_i\mid i\in[n]\}$, obtained by converting a size-$m$ $k$-CSP as given by Definition6. For convenience, we will index the polynomial equations $E_j$ in binary notation, so for example if $b_1\ldots b_{\log m}\in\{0,1\}^{\log m}$ is the binary notation for $j\in[m]$, we will write $E_j$ as $E_{b_1\ldots b_{\log m}}$. We define a new set of polynomial equations $F^{\mathsf{SEL}}_{n,k}$ with parameters $m,m'$ as follows. The variables are $\mathbf x\cup\mathbf y$, where $\mathbf x$ are the original variables of $F_{n,k}$ and $\mathbf y=\{y_{i,j},\bar y_{i,j}\mid i\in[m'],j\in[\log m]\}$ are new ``pigeon'' variables. Let $E^{\mathsf{SEL}}$ be the following set of equations, where $\mathbf y_i\neq b_1\ldots b_{\log m}$ abbreviates the monomial $\prod_{b_j=1}y_{i,j}\prod_{b_j=0}\bar y_{i,j}$:
\begin{enumerate}[label=(\roman*)]
\item $\forall i\in[m'],\ \forall b_1\ldots b_{\log m}\in\{0,1\}^{\log m},\ (\mathbf y_i\neq b_1\ldots b_{\log m})\cdot E_{b_1\ldots b_{\log m}}=0$;
\item $\forall i,i'\in[m'],i\neq i',\ \forall b_1\ldots b_{\log m}\in\{0,1\}^{\log m},$
\[(\mathbf y_i\neq b_1\ldots b_{\log m})\cdot(\mathbf y_{i'}\neq b_1\ldots b_{\log m})=0.\]
\end{enumerate}
$F^{\mathsf{SEL}}_{n,k}$ consists of the polynomial equations $E^{\mathsf{SEL}}$ together with the Boolean axioms $B_{i,j}=\{y_{i,j}^2-y_{i,j}=0,\bar y_{i,j}^2-\bar y_{i,j}=0,y_{i,j}\bar y_{i,j}=0\}$ for all $i\in[m'],j\in[m]$.
Intuitively we think of the $y$ variables as a mapping from $m'$ pigeons to $m$ holes, where the holes correspond to the $m$ axioms/constraints from $E$. For $i\in[m']$, the $i$th ``pigeon'' $\mathbf y_i$ selects a hole (an equation from $E$).
The first set of polynomial equations in $E^{\mathsf{SEL}}$ states that if pigeon $\mathbf y_i$ selects the equation $E_{b_1\ldots b_{\log m}}$, then this equation must be satisfied; the second set of equations in $E^{\mathsf{SEL}}$ states that the mapping is one-to-one and thus altogether the $y$ selector variables choose a subset $E'$ of exactly $m'$ equations from $E$. Thus $F^{\mathsf{SEL}}_{n,k}$ asserts that there exists a subset of $m'$ constraints of $F_{n,k}$ (chosen by the $y$-variables) that are satisfiable.
Throughout this paper, the $x$-variables are the variables that underly $F_{n,k}$; the $y$-variables are the selector/pigeon variables described above that choose a subset of $m'$ constraints from $F_{n,k}$, and the extension variables used in the PC+Ext refutation will be the $z$-variables.
Our hard instances will be $F^{\mathsf{SEL}}_{n,k}$, with $m=10n$, $m'=(1-\epsilon)m$, where $F_{n,k}$ is (the polynomial translation of) an unsatisfiable $k$-CSP formula with $m=10n$ $k$-local constraints over variables $\mathbf x=x_1\ldots x_n$, satisfying the follow property:
\paragraph{Property8.} Every subset of $(1-\epsilon)m'$ constraints is unsatisfiable and requires linear PC degree.
The following Theorem shows that for sufficiently large $n$, such formulas exist. Similar proofs have appeared in several papers (e.g., [5]) but we give a proof in the Appendix for completeness.
\paragraph{Theorem9.} Let $m=10n$. Then there exists constants $k>0$, $0<\epsilon<1$ such that for sufficiently large $n$, there exists $k$-CSP formulas $\{F_{n,k}\}$ with $m$ constraints such that Property8 holds with $m'=(1-\epsilon)m$.

\subsection*{4.2 Singular and Nonsingular variables}
Let us fix the finite field $\mathbb F_p$, $p>2$ for the rest of the article. We also fix a set of unsatisfiable polynomials $F$ over Boolean variables $\mathbf x\cup\mathbf y$, and a set of extension axioms $\mathsf{Ext}$ of the form $z-Q$ over variables $\mathbf z$. Whenever we refer to a refutation $\Pi$, we assume that it is a PC+Ext refutation of $F\cup\mathsf{Ext}$.
\paragraph{Definition10 (Support of a variable).}
Let $z-Q(w_{i_1},\ldots,w_{i_{\kappa}})=0$ be a $\kappa$-local extension axiom associated with $z$. We define the set $\vars(Q)=\{w_{i_1},\ldots,w_{i_{\kappa}}\}$ and sometimes write $\vars(z)$ to denote $\vars(Q)$, the set of variables that $z$ depends on. The support of $z$, $\supp(z)\subseteq[0,p-1]$, is equal to the set of all values $a\in[0,p-1]$ such that there exists a Boolean assignment $\alpha$ to the variables of $Q$ such that $Q(\alpha)=a$. Sometimes we also indicate this by $\supp(Q)$.
We extend the definition of support also to Boolean variables. For a Boolean variable $w$, $\supp(w)=\{0,1\}$ as enforced by the Boolean axiom $w^2=w$.
\paragraph{Definition11 (Singular and Nonsingular variables w.r.t. Ext).}
Let $\mathsf{Ext}$ be a set of extension axioms and let $z$ be an extension variable with an axiom in $\mathsf{Ext}$. We say that $z$ is Singular w.r.t. $\mathsf{Ext}$ iff $0\in\supp(z)$; otherwise we say that $z$ is Nonsingular w.r.t. $\mathsf{Ext}$. Any Boolean variable is considered Singular by default, independent of the set $\mathsf{Ext}$, since zero always belongs to its support. For a term $t$, let $\sing(t)$ be the subterm of $t$ containing the Singular variables in $t$, and let $\nsing(t)$ be the subterm of $t$ containing the Nonsingular variables.
Note that for a Singular extension variable $z$, it is possible to set $z$ to zero, However, we note that this may falsify other polynomial equations in $F$. For example, if $xy=0$ is a polynomial equation in $F$, and the extension axiom for $z$ is $z-1+xy=0$, then setting $x=y=1$ forces $z=0$, but this falsifies $xy=0$.
\paragraph{Definition12.} Let $A\subseteq[1,\ldots,p-1]$, $A\neq\emptyset$. Define $\ell(A)$ to be the least $\ell\in[1,p-1]$ such that the set $\{a^{\ell}\mid a\in A\}$ is singleton. For a Nonsingular $z$, define $\ell(z)=\ell(\supp(z))$.
\paragraph{Lemma13.} Let $z$ be a Nonsingular extension variable with extension axiom $z-Q=0$. Then the following polynomial equations are implied by (and therefore derivable from) the extension axiom for $z$ plus the Boolean axioms for all variables in $\vars(Q)$, in degree at most $|\vars(Q)|$.
\begin{enumerate}
\item $z-Q'=0$, where $Q'$ is the multilinear version of $Q$;
\item For any $A'\subseteq[0,p-1]$ such that $\supp(z)\subseteq A'$, $\prod_{a\in A'}(z-a)=0$;
\item $z^{\ell(z)}-c=0$ for some $c\in\mathbb F_p^*$.
\end{enumerate}
In particular, if $z$ is Nonsingular, then the polynomial equation $z^{p-1}-1=0$ is implied by $z-Q=0$ together with the Boolean axioms for $\vars(Q)$.
\begin{proof}
Let $z-Q(w_{i_1},\ldots,w_{i_{\kappa}})=0$ be the extension axiom for $z$, and let $\supp(z)=A\subseteq A'\subseteq[1,p-1]$. First, we can derive the multilinear version of $Q$, $Q'$, from $Q$ together with the Boolean axoms $w^2-w=0$ for all $w\in\vars(Q)$. Secondly, by definition, $\supp(z)=A$ means that the allowable values for $z$ over Boolean assignments to $\vars(Q)$ are the values in $A$. Therefore, $z-Q=0$ together with the Boolean axioms $w^2-w=0$ for all $w\in\vars(Q)$ implies $\prod_{a\in A}(z-a)=0$. Furthermore, this polynomial has a PC derivation, by the derivational completeness of PC. Since $A\subseteq A'$, $\prod_{a\in A'}(z-a)=0$ is a weakening of $\prod_{a\in A}(z-a)=0$ and is therefore derivable from $\prod_{a\in A}(z=a)=0$. Lastly, we will argue that there exists some constant $c\in\mathbb F_p^*$ such that $z^{\ell(A)}-c=0$ is semantically implied by $z-Q=0$ plus the Boolean axioms for $\vars(z)$ and therefore is derivable from these axioms. Since the only allowable values for $z$ under the Boolean axioms are the values in $A$, and since by definition of $\ell(A)$, for every $a\in A$, $a^{\ell(A)}=c$ for some $c\in\mathbb F_p^*$, it follows that $z^{\ell(A)}-c=0$.
\end{proof}
\paragraph{Definition14.} For a term $t$ and a variable $w$, $\deg(t,w)$ is equal to the degree of $w$ in $t$. If $w$ is Nonsingular, then $w^{p-1}=1\bmod p$, so $\deg(t,w)<p-1$. On the other hand if $w$ is Singular then we have $w^p=w\bmod p$ and therefore $\deg(t,w)<p$. For a term $t$ the degree of $t$, $\deg(t)$, equals $\sum_{w\in\vars(t)}\deg(t,w)$.
\subsection*{4.3 Quadratic degree}
The next definition is a generalization/modification of Sokolov's definition of Quadratic degree for the more general scenario where the proof contains extension variables that are Singular as well as ones that are Nonsingular.
\paragraph{Definition15 (Quadratic degree).}
Let $V$ be a set of variables and let $S$ be a subset of $V$. For a pair of terms $t_1,t_2$ over $V$, and a variable $w\in V$, we define $\Qdeg_S(t_1,t_2,w)$ as follows. If $w\in S$, then $\Qdeg_S(t_1,t_2,w)=1$ if $w$ occurs in at least one of $t_1$ or $t_2$; if $w\notin S$, then $\Qdeg_S(t_1,t_2,w)=1$ if and only if $\deg(t_1,w)\neq\deg(t_2,w)$. The overall quadratic degree of the pair $t_1,t_2$, $\Qdeg_S(t_1,t_2)$, is equal to $\sum_{w\in V}\Qdeg_S(t_1,t_2,w)$. The quadratic degree of a polynomial $P$ is equal to the maximum quadratic degree over all pairs $(t_1,t_2)$ such that $t_1,t_2\in P$. For a proof $\Pi$, the quadratic degree of $\Pi$ is the maximum quadratic degree over all polynomials $P\in\Pi$.
We usually instantiate the above definition with $V=\mathbf x\cup\mathbf y\cup\mathbf z$ and with $S$ being the set of Singular variables as defined by the extension axioms corresponding to $\mathbf z$. However since $\Qdeg_S$ is a different measure for every $S$, and our set of Singular variables can change under the application of a restriction $\rho$ to the variables in $V$, we must make sure that our measure of Quadratic degree does not change significantly under a restriction.\footnote{If the set $S$ does not change under a restriction, $\Qdeg_S$ can still change under the restriction as terms can shrink or disappear when variables are set by the restriction. However, this is no different from how the usual notion of degree changes under a restriction, and it is trivial to show that $\Qdeg_S$ always decreases. Therefore we ignore this for the sake of simplicity.} Fortunately, we can show that for any two sets $S$ and $T$ such that $T\subseteq S$, $\Qdeg_T\leq\Qdeg_S$. Along with the simple observation that the set of Singular extension variables can only shrink under a restriction, this implies that our measure of Quadratic degree can only decrease under a restriction. We make this formal below.
\paragraph{Lemma16.} Let $V$ be a set of variables and let $S$ and $T$ be subsets of $V$ such that $T\subseteq S$. Then for any two terms $t_1,t_2$ over $V$, $\Qdeg_T(t_1,t_2)\leq\Qdeg_S(t_1,t_2)$.
\begin{proof}
Note that for a variable $w\in S-T$, $\Qdeg_S(t_1,t_2,w)=1$ when $w$ has a nonzero exponent in one of $t_1$ or $t_2$, otherwise zero. However, $\Qdeg_T(t_1,t_2,w)=1$ if and only if the previous condition is satisfied and the exponents of $w$ in $t_1$ and $t_2$ are not equal. Thus the claim follows.
\end{proof}
Henceforth, when we refer to Quadratic degree, we always fix the set $S$ to be the set of Singular variables w.r.t. the underlying extension axioms. We have the following important corollary that this measure always decreases under a restriction to the underlying variables.
\paragraph{Corollary17.} Let $F$ be a set of unsatisfiable polynomials over variables $\mathbf x\cup\mathbf y$ and let $\mathsf{Ext}$ be a set of extension axioms of the form $z-Q(w_{i_1},\ldots,w_{i_{\kappa}})$ for variables $z\in\mathbf z$ and $w_{i_1},\ldots,w_{i_{\kappa}}\in\mathbf x\cup\mathbf y$. Let $\rho$ be a restriction to $\mathbf x\cup\mathbf y$ and let $\mathsf{Ext}|_{\rho}$ be the axioms given by $z-Q|_{\rho}$ for each axiom $z-Q\in\mathsf{Ext}$. The Quadratic degree w.r.t. $\mathsf{Ext}|_{\rho}$ is at most the Quadratic degree w.r.t. $\mathsf{Ext}$.
\begin{proof}
Since $\supp(Q|_{\rho})\subseteq\supp(Q)$ for any polynomial $Q$, we have that the set of Singular variables under $\mathsf{Ext}|_{\rho}$ is a subset of those under $\mathsf{Ext}$. Therefore our claim follows from the previous lemma.
\end{proof}
\paragraph{Lemma18 (Quadratic degree upper bounds degree of Singular variables).} For any term $t$, $\deg(\sing(t))\leq p\Qdeg(t,t)$.
\begin{proof}
For any Singular variable $w$, $\Qdeg(t,t,w)=1$ if and only if $w$ occurs in $t$. Since $w$ can occur in $t$ with degree at most $p-1$, the claim follows.
\end{proof}
\paragraph{Definition19 (High quadratic degree terms).} For a proof $\Pi$ and $d\geq0$, let $H_d(\Pi)$ denote the set of unordered pairs $(t_1,t_2)$ of quadratic degree at least $d$. That is, $H_d(\Pi)$ is the set of unordered pairs of terms $(t_1,t_2)$ such that $t_1,t_2$ both occur in $P$ for some polynomial $P\in\Pi$, and $\Qdeg(t_1,t_2)\geq d$.
\paragraph{Lemma20.} Let $\Pi$ be a PC+Ext refutation of $F$ and let $z$ be a Nonsingular variable. Let $\Pi'$ be the proof obtained from $\Pi$ by reducing each line of $\Pi$ by $z^{\ell(z)}-c=0$ for some $c\in\mathbb F_p^*$. Then $|H_d(\Pi')|\leq|H_d(\Pi)|$ for any $d\geq0$.
\begin{proof}
Consider a polynomial $P\in\Pi$ and a pair of terms $(t_1,t_2)$ that occur in $P$. For any variable $w$ distinct from $z$, $\Qdeg(t_1,t_2,w)$ is unaltered when $P$ is reduced by $z^{\ell(z)}=c$. On the other hand, if $z$ does not contribute to the Quadratic degree of $(t_1,t_2)$ i.e. $\Qdeg(t_1,t_2,z)=0$, then it will still be $0$ after reducing by $z^{\ell(z)}=c$. Therefore $\Qdeg(t_1,t_2)$ never increases for any pair $(t_1,t_2)$ and thus $|H_d(\Pi')|\leq|H_d(\Pi)|$.
\end{proof}
The following is a generalized version of the argument from [20] that shows how to convert a proof with low Quadratic degree to one with low degree.
\paragraph{Lemma21.} Let $F$ be a set of unsatisfiable polynomials of degree $d_0$ with a PC refutation of Quadratic degree at most $d\geq d_0$ over $\mathbb F_p$. Then $F$ has a PC refutation of degree at most $3pd$.
\begin{proof}
The proof of this lemma is largely based on (a slightly cleaner version of) Sokolov's argument ([20], Lemma3.6) that low Quadratic degree over $\{\pm1\}$ variables implies low degree. Our first observation is that Sokolov's argument can be applied to any refutation of low Quadratic degree over $\mathbb F_p$ such that every term contains only Nonsingular variables. In particular if $\{P_j\}$ is a refutation that only contains Nonsingular terms, then we can use his argument to show that $\{t_j^{p-2}P_j\}$ is also a valid refutation for some carefully chosen term $t_j\in P_j$. Moreover, the degree of the latter refutation is bounded by a constant times the Quadratic degree of the former one. To see this, first note that for two Nonsingular terms $t_1$ and $t_2$, we have that $\deg(t_1t_2^{p-2})\leq(p-1)\cdot\Qdeg(t_1,t_2)$, because of the following. For a variable $z$ that is Nonsingular such that $z$ occurs in $t_1$ and $t_2$ with $\deg(t_1,z)=\deg(t_2,z)$, we have $\deg(t_1t_2^{p-2},z)=\Qdeg(t_1,t_2,z)=0$ since it would appear in $t_1t_2^{p-2}$ with an exponent that is a multiple of $p-1$, and $z^{p-1}=1$ holds for Nonsingular variables. Any other Nonsingular $z$ that occurs in at least one of $t_1$ and $t_2$ has $\deg(t_1t_2^{p-2},z)<p-1$ and $\Qdeg(t_1,t_2,z)=1$. Therefore the degree of $t_1t_2^{p-2}$ is at most $p\cdot\Qdeg(t_1,t_2)$ when $t_1$ and $t_2$ contain only Nonsingular variables. This implies that the lines in the new refutation $\{t_j^{p-2}P_j\}$ have degree at most $p$ times the Quadratic degree of the original refutation $\{P_j\}$. Sokolov additionally showed that each line in the new refutation can be derived from previous lines without exceeding degree equal to $2p$ times the Quadratic degree of the original refutation, completing the argument.
In our case we deal with terms containing both Singular and Nonsingular variables. The above argument cannot be applied directly to our case, since it crucially depends on the fact that Nonsingular variables can be raised to the power $p-1$ to make them vanish. Fortunately by Lemma18, the degree of Singular variables in any term is at most $p$ times the Quadratic degree with itself. Given this bound, we can ignore for each term the part that contains Singular variables, and apply the above argument only with respect to the Nonsingular part of each term, to reduce the degree of Nonsingular variables in each term of the refutation. Since we now have a bound on the degree of both Singular and Nonsingular variables in each term, we have bounded its degree. We describe this in full technical detail below.
Let $\{P_j\}$ be a refutation of $F$ with Quadratic degree bounded by $d$. For any term $t$ recall that $\nsing(t)$ denotes the subterm of $t$ containing only Nonsingular variables. Note that $\nsing(t)^{p-1}=1$ for any $t$. For every line $P_j$ in the refutation, we pick a term $t_j\in P_j$ and define $P'_j=\nsing(t_j)^{p-2}P_j$. Note that by the arguments outlined above, for any two terms $t_1$ and $t_2$ in $P_j$, we have $\deg(\nsing(t_1)^{p-2}\nsing(t_2))\leq pd$ and thus the degree of Nonsingular variables in any term of $P'_j$ is bounded by $pd$. Since the Singular variables in any term remain unchanged under multiplication by $\nsing(t_j)^{p-2}$, the Singular degree of $P'_j$ the same as that of $P_j$ and is bounded by $pd$ (Lemma18) and therefore $\deg(P'_j)\leq pd+pd=2pd$. We now show that the set $\{P'_j\}$ forms a valid refutation of $F$ and each $P'_j$ can be derived from previous lines in degree $3pd$. If $P_j$ is one of the axioms, we multiply by $\nsing(t_j)^{p-2}$ to get $P'_j$ for an arbitrary $t_j\in P_j$, and this takes degree $pd_0\leq pd$.
If $P_j=wP_{j_1}$ for $j_1<j$ and some variable $w$, we choose $t_j\in P_j$ such that $t_j=wt_{j_1}$ where $t_{j_1}\in P_{j_1}$ was chosen earlier. If $w$ is Singular, we have $\nsing(t_j)=\nsing(t_{j_1})$ and therefore $P'_j=\nsing(t_j)^{p-2}P_j=w\cdot\nsing(t_{j_1})^{p-2}P_{j_1}=wP'_{j_1}$. On the other hand, if $w$ is Nonsingular, we have $\nsing(t_j)=w\cdot\nsing(t_{j_1})$ and therefore $P'_j=\nsing(t_j)^{p-2}P_j=w^{p-1}\cdot\nsing(t_{j_1})^{p-2}P_{j_1}=P'_{j_1}$. Finally, let $P_j=P_{j_1}+P_{j_2}$ for $j_1,j_2<j$. We pick an arbitrary term $t_j\in P_j$. Note that since $\nsing(t)^{p-1}=1$ for any term $t$, $P_{j_1}=\nsing(t_{j_1})P'_{j_1}$ and $P_{j_2}=\nsing(t_{j_2})P'_{j_2}$ and thus we have
\begin{multline*}
P'_j=\nsing(t_j)^{p-2}\nsing(t_{j_1})P'_{j_1}\\
+\nsing(t_j)^{p-2}\nsing(t_{j_2})P'_{j_2}
\end{multline*}
for $t_{j_1}\in P_{j_1}$ and $t_{j_2}\in P_{j_2}$ chosen earlier. We now show that $\deg(\nsing(t_j)^{p-2}\nsing(t_{j_1}))\leq pd$ and $\deg(\nsing(t_j)^{p-2}\nsing(t_{j_2}))\leq pd$ to conclude the proof. Since every term in $P_j$ appears in one of $P_{j_1},P_{j_2}$, let $t_j\in P_{j_1}$ without loss of generality. Then we have that $t_j,t_{j_1}$ both appear in $P_{j_1}$ and thus $\deg(\nsing(t_j)^{p-2}\nsing(t_{j_1}))\leq pd$. If $t_{j_2}\in P_j$ i.e. it is not cancelled in the sum $P_{j_1}+P_{j_2}$, then we have $t_j,t_{j_2}$ both appear in $P_j$ and hence $\deg(\nsing(t_j)^{p-2}\nsing(t_{j_2}))\leq pd$. If $t_{j_2}\notin P_j$, this implies that it was cancelled in the sum $P_{j_1}+P_{j_2}$ and therefore $t_{j_2}\in P_{j_1}$ and $\deg(\nsing(t_j)^{p-2}\nsing(t_{j_2}))\leq pd$.
\end{proof}

\subsection*{4.4 The Split Operation}
In this section we will show how to apply a restriction and then use an operation Split (motivated by [20]) in order to eliminate high quadratic degree terms. Our main focus will be to handle the case where the variable to be set is an extension variable with extension axiom $z-Q=0$ where $z$ is Nonsingular, since in the other case we can potentially just set $z=0$ to eliminate terms. We start by showing how to apply a small Boolean restriction $\rho$ such that $Q|_{\rho}$ is a simple linear function of just one variable.
\paragraph{Lemma22.} Let $z$ be an extension variable with extension axiom $z-Q(w_1,\ldots,w_k)$, for $k\leq\kappa$. Assume that $z$ is Nonsingular (i.e. $\supp(Q)\subseteq[1,\ldots,p-1]$) and $|\supp(Q)|\geq2$. Then for every $l\in[0,\ldots,\ell(\supp(Q))-1]$, there exists a variable $w^*$ in $\vars(Q)$, and a restriction $\delta$ to $\vars(Q)-w^*$ such that:
\begin{enumerate}[label=(\arabic*)]
\item $Q|_{\delta}=(b-a)w^*+a$, where $b,a\in\supp(Q)$. Thus $Q|_{\delta}$ is a linear function of $w^*$ and $\supp(Q|_{\delta})=\{a,b\}$;
\item $a^l\neq b^l\pmod p$.
\end{enumerate}
\begin{proof}
We will create a decision tree that will query $\vars(Q)$ one-by-one. Associated with the root $r$ is the set of values $S_r=\{a^l\mid a\in\supp(Q)\}$. That is, we label the root with the set of all possible values that $z^l$ can take on. Since $l<\ell(\supp(z))$, it follows that $|S_r|\geq2$ (since otherwise we would have $l=\ell(\supp(z))$). At the root we query the first variable $w_1$, with left edge labelled by $w_1=0$ and right edge labelled by $w_1=1$. Now we label the left vertex with the set $\{a^l\mid a\in\supp(Q|_{w_1=0})\}$, of all values that $z^l$ can take on under the restriction $w_1=0$. Similarly we label the right vertex with the set $\{a^l\mid a\in\supp(Q|_{w_1=1})\}$. We continue recursively, querying the next variable at each vertex $v$ of the decision tree, as long as the set of allowable values for $z^l$ under the partial restriction $\rho_v$ associated with $v$ is greater than one. Now consider the longest path, $\xi$ in $T$. The partial restriction $\rho$ associated with $\xi$ sets the first $k'$ variables, where $k'\geq1$ since initially $z^l$ takes on at least two values. Also since $\xi$ is a complete path, the associated set $\{a^l\mid a\in\supp(Q|_{\rho})\}$ contains exactly one element, call it $q$.
Now consider the twin path $\xi'$ with associated restriction $\rho'$, where $\rho'$ is obtained from $\rho$ by toggling the value of the last variable, $w_{k'}$, queried. Again since $\xi'$ is a complete path, the associated set $\{a^l\mid a\in\supp(Q|_{\rho'})\}$ contains exactly one element, call it $q'$. Note that $q,q'$ must be distinct.
Let $\delta$ be the following assignment to $\vars(Q)-w_{k'}$: for $1\leq j<k'$, we set $\delta(w_j)=\rho(w_j)=\rho'(w_j)$, and for $k'<j\leq k$, we set $\delta(w_j)=0$. Setting $w^*=w_{k'}$, $Q|_{\delta}$ is a linear equation of the form $(b-a)w^*+a$, where $b,a\in\supp(Q)$. Finally, by construction, $a^l\neq b^l$ (since otherwise the two paths corresponding to $\rho,\rho'$ would be the same).
\end{proof}
In the remainder of this subsection, we will be interested in the case where we want to eliminate some Nonsingular extension variable $z$ from the refutation, and we have already applied the above Lemma so that the extension axiom for $z$ is of the form $z-((b-a)w+a)=0$, where $w$ is some variable in $\mathbf x\cup\mathbf y$. Thus, $\supp(z)=\{a,b\}$. The next two Lemmas generalizes a similar argument due to Sokolov, and show how to remove Quadratic degree pairs of the form $(t_1z^i,t_2z^j)$ for a carefully chosen pair $i,j$ from the refutation via the Split operation.
\paragraph{Lemma23.} Let $z$ be an extension variable such that $\supp(z)=\{a,b\}$, where $a\neq b$ and $a,b\in\mathbb F_p^*$ and let $P$ be any polynomial. Then, for any two distinct numbers $\ell_0,\ell_1$ where $\ell_0<\ell_1$ and $a^{\ell_1-\ell_0}\neq b^{\ell_1-\ell_0}$, there exists a unique polynomial $R=R_0z^{\ell_0}+R_1z^{\ell_1}$ such that $R=P\bmod(z-a)(z-b)$. That is, $R(a)=P(a)$ and $R(b)=P(b)$, where $P(a)$ denotes the polynomial $P$ under the substitution $z=a$.
\begin{proof}
Let $z-Q=0$ be the extension axiom for $z$, where $\supp(z)=\{a,b\}$. Then by Lemma13 the polynomial $(z-a)(z-b)=0$ is implied by (and derivable from) the extension axiom for $z$ plus the Boolean axioms. We can assume without loss of generality that $P$ has the form $P_0+zP_1+\ldots+z^{p-2}P_{p-2}$.
Now we want to argue that there exists a polynomial $R=z^{\ell_0}R_0+z^{\ell_1}R_1$, where $R_0,R_1$ are polynomials over $\vars(P)-z$, and such that $R(a)=P(a)$, and $R(b)=P(b)$. We can find $R_0$ and $R_1$ by solving the following system of equations, where we view $R_0,R_1$ as the underlying variables, and treating $P(a),P(b)$ as constants:
\[a^{\ell_0}R_0+a^{\ell_1}R_1=P(a),\qquad b^{\ell_0}R_0+b^{\ell_1}R_1=P(b).\]
This has a (unique) solution since the determinant of the associated matrix is
\[\begin{vmatrix}a^{\ell_0}&a^{\ell_1}\\b^{\ell_0}&b^{\ell_1}\end{vmatrix}=a^{\ell_0}b^{\ell_0}(b^{\ell_1-\ell_0}-a^{\ell_1-\ell_0}).\]
By our assumption, this matrix is non-singular over $\mathbb F_p$ and therefore the above system of equations has a unique solution over $\mathbb F_p$, given by:
\[\binom{R_0}{R_1}=\begin{pmatrix}a^{\ell_0}&a^{\ell_1}\\b^{\ell_0}&b^{\ell_1}\end{pmatrix}^{-1}\binom{P(a)}{P(b)}.\]
Abbreviating $a^{\ell_0},a^{\ell_1},b^{\ell_0},b^{\ell_1}$ by $a_0,a_1,b_0,b_1$ respectively, we have by definition of the inverse:
\[\binom{R_0}{R_1}=\begin{pmatrix}a_0&a_1\\b_0&b_1\end{pmatrix}^{-1}\binom{P(a)}{P(b)}=\frac{1}{a_0b_1-a_1b_0}\begin{pmatrix}b_1&-a_1\\-b_0&a_0\end{pmatrix}\binom{P(a)}{P(b)}.\]
Solving for $R_0$ we have:
\begin{align*}
R_0&=\frac{b_1}{a_0b_1-a_1b_0}P(a)-\frac{a_1}{a_0b_1-a_1b_0}P(b)\\
&=\frac{b_1}{a_0b_1-a_1b_0}\left(a_0P_{\ell_0}+a_1P_{\ell_1}+\sum_{i\neq\ell_0,\ell_1}a^iP_i\right)\\
&\quad-\frac{a_1}{a_0b_1-a_1b_0}\left(b_0P_{\ell_0}+b_1P_{\ell_1}+\sum_{i\neq\ell_0,\ell_1}b^iP_i\right)\\
&=\frac{a_0b_1}{a_0b_1-a_1b_0}P_{\ell_0}+\frac{a_1b_1}{a_0b_1-a_1b_0}P_{\ell_1}\\
&\quad+\frac{b_1}{a_0b_1-a_1b_0}\left(\sum_{i\neq\ell_0,\ell_1}a^iP_i\right)-\frac{a_1b_0}{a_0b_1-a_1b_0}P_{\ell_0}\\
&\quad-\frac{a_1b_1}{a_0b_1-a_1b_0}P_{\ell_1}-\frac{b_1}{a_0b_1-a_1b_0}\left(\sum_{i\neq\ell_0,\ell_1}b^iP_i\right)\\
&=P_{\ell_0}+\sum_{i\neq\ell_0,\ell_1}c_{0i}P_i
\end{align*}
for some constants $c_{0i}\in\mathbb F_p$. And similarly solving for $R_1$, it has the following form:
\[R_1=P_{\ell_1}+\sum_{i\neq\ell_0,\ell_1}c_{1i}P_i\]
for some constants $c_{1i}\in\mathbb F_p$.
\end{proof}
\paragraph{Definition24 (Split).} Let $z$ be an extension variable with extension axiom $z-Q=0$ such that $\supp(z)=\{a,b\}\subseteq[1,\ldots,p-1]$. For any polynomial $P$ and for every $\ell_0<\ell_1$ such that $a^{\ell_1-\ell_0}\neq b^{\ell_1-\ell_0}$, let $R=R_0z^{\ell_0}+R_1z^{\ell_1}$ be the unique polynomial given by Lemma23 such that $R=P\bmod(z-a)(z-b)$. Then $\Split_{z,\ell_1,\ell_0}(P)$ is defined to be the pair of polynomials $\{R_0,R_1\}$. For a proof $\Pi$, and an extension variable $z$ such that $\supp(z)=\{a,b\}$, we define $\Split_{z,\ell_0,\ell_1}(\Pi)$ to be the sequence of lines $\Split_{z,\ell_0,\ell_1}(P)$, over all $P\in\Pi$.
\paragraph{Lemma25.} Let $\Pi$ be a refutation of a set of unsatisfiable polynomials $F$. Let $z$ be a variable that occurs in $\Pi$ such that the polynomials in $F$ do not contain $z$ except for the axiom $(z-a)(z-b)=0$ for some $a,b\in\mathbb F_p^*$. Then for any $\ell_0,\ell_1$ such that $\ell_0<\ell_1$ and $a^{\ell_1-\ell_0}\neq b^{\ell_1-\ell_0}$, $\Pi'=\Split_{z,\ell_0,\ell_1}(\Pi)$ forms a valid refutation of $F$ modulo $(z-a)(z-b)$.
\begin{proof}
Fix an extension variable $z$ in $\Pi$ such that it does not occur in any axioms except $(z-a)(z-b)=0$, and let $\ell_0,\ell_1$ be such that $\ell_0<\ell_1$ and $a^{\ell_1-\ell_0}\neq b^{\ell_1-\ell_0}$. We will show by induction on the number of lines in $\Pi$ that $\Split_{z,\ell_0,\ell_1}(\Pi)$ is a valid derivation that meets the conditions of the lemma. For the base case, note that all of the axioms are either free of $z$ or eliminated as a result of reducing by $(z-a)(z-b)$, and hence their Split versions are derivable. Now suppose that the Lemma holds for the first $j-1$ lines of $\Pi$; that is, $\Split_{z,\ell_0,\ell_1}(\Pi_{j-1})$ is a valid derivation, where $\Pi_{j-1}$ denotes the first $j-1$ lines of $\Pi$.
The first case is where $P_j$ is a linear combination of two previously derived lines, so $P_j=\alpha P_{j_1}+\beta P_{j_2}$ for some $j_1$ and $j_2$ less than $j$ and $\alpha,\beta\in\mathbb F_p$. Using the inductive hypothesis, we have:
\begin{align*}
P_j&=\alpha(z^{\ell_0}R_{j_10}+z^{\ell_1}R_{j_11})+\beta(z^{\ell_0}R_{j_20}+z^{\ell_1}R_{j_21})\pmod{(z-a)(z-b)}\\
&=z^{\ell_0}(\alpha R_{j_10}+\beta R_{j_20})+z^{\ell_1}(\alpha R_{j_11}+\beta R_{j_21})\pmod{(z-a)(z-b)}.
\end{align*}
By the uniqueness of the polynomial $R_j=z^{\ell_0}R_{j0}+z^{\ell_1}R_{j1}$ that is equivalent to $P_j\bmod(z-a)(z-b)$ (by Lemma23), this implies that $R_{j0}=\alpha R_{j_10}+\beta R_{j_20}$ and similarly $R_{j1}=\alpha R_{j_11}+\beta R_{j_21}$, and thus $R_{j0}$ can be derived from a linear combination of $R_{j_10}$ and $R_{j_20}$ and similarly for $R_{j1}$.
The second case is when $P_j$ is derived from a previously derived line $P_{j'}$ by multiplying $P_{j'}$ by a variable $w$. That is, $P_j=wP_{j'}$ for some $j'<j$. If $w\neq z$, then we have that $R_{j1}=wR_{j'1}$ (similarly for $R_{j0}$). If $w=z$ then we have:
\[\binom{R_{j'1}}{R_{j'0}}=\begin{pmatrix}a^{\ell_1}&a^{\ell_0}\\b^{\ell_1}&b^{\ell_0}\end{pmatrix}^{-1}\binom{P_{j'}(a)}{P_{j'}(b)}\]
from which we need to derive
\begin{align*}
\binom{R_{j1}}{R_{j0}}&=\begin{pmatrix}a^{\ell_1}&a^{\ell_0}\\b^{\ell_1}&b^{\ell_0}\end{pmatrix}^{-1}\binom{P_j(a)}{P_j(b)}\\
&=\begin{pmatrix}a^{\ell_1}&a^{\ell_0}\\b^{\ell_1}&b^{\ell_0}\end{pmatrix}^{-1}\binom{aP_{j'}(a)}{bP_{j'}(b)}\\
&=\begin{pmatrix}a^{\ell_1}&a^{\ell_0}\\b^{\ell_1}&b^{\ell_0}\end{pmatrix}^{-1}\begin{pmatrix}a&0\\0&b\end{pmatrix}\begin{pmatrix}a^{\ell_1}&a^{\ell_0}\\b^{\ell_1}&b^{\ell_0}\end{pmatrix}\binom{R_{j'1}}{R_{j'0}}.
\end{align*}
Thus, $R_{j1}$ can be derived as a linear combination of $R_{j'1}$ and $R_{j'0}$, and similarly for $R_{j0}$.
\end{proof}

\subsection*{4.5 Proof of Main Theorem}
The proof of our lower bound for the tautology $F^{\mathsf{SEL}}_{n,k}$ with extension axioms $\mathsf{Ext}$ proceeds by choosing a variable in the given refutation $\Pi$ that contributes to a lot of high quadratic degree pairs of terms in $\Pi$. If this variable is Singular, we apply the restriction that sets it to zero. On the other hand, if it is Nonsingular and therefore an extension variable $z$, we first reduce it to depend on a single variable $w^*$ by applying a restriction chosen from Lemma22, and then use a more complicated case analysis (see Lemma30) in order to apply the Split operation from Lemmas23 and25 on $z$. In both of these cases we are able to remove a small fraction of high Quadratic degree terms, and thus after sufficiently many iterations we obtain a refutation of low Quadratic degree. We convert this to a refutation of low (usual) degree using Lemma21, and then substitute for the pigeon variables $\mathbf y$ to select a subset of equations from $F_{n,k}$ that require high degree, obtaining a contradiction.
\subsubsection*{4.5.1 Cleanup operations}
In order to get the contradiction at the end of the above argument, we need to ensure that our process above is always working with a subset of equations of $F_{n,k}$ that are untouched, i.e. unaffected by earlier restrictions to variables. We also need to eliminate any partially assigned pigeons so that we have full choice over the equations we are able to pick at the end. Additionally, a key requirement of the Split lemmas (Lemmas23 and25) is that the variable $z$ we Split on must not appear in any axioms except for one of the form $(z-a)(z-b)=0$, which indicates that it takes two distinct values. In particular, we cannot set $z$ or the underlying variable $w^*$ described above in order to eliminate them from the refutation. This presents us with a unique requirement: for any choice of a variable $w^*\in\mathbf x\cup\mathbf y$, we need to be able to eliminate all axioms containing $w^*$ without actually setting it. We show how to perform these operations by making use of the structure of our tautology $F^{\mathsf{SEL}}_{n,k}$.
We first show how to ``ban'' an equation $E_{b_1\ldots b_{\log m}}$ from $F_{n,k}$ by switching to a set of axioms that prevent any pigeon from being assigned to $b_1\ldots b_{\log m}$.
\paragraph{Lemma26.} Let $\Pi$ be a refutation of $F^{\mathsf{SEL}}_{n,k}|_{\rho}$ for some restriction $\rho$ and let $(\mathbf y_i\neq b_1\ldots b_{\log m})\cdot E_{b_1\ldots b_{\log m}}=0$ be one of its axioms. Then there exists another valid refutation $\Pi'$ with the latter axiom replaced by the axiom $(\mathbf y_i\neq b_1\ldots b_{\log m})\equiv\prod_{b_j=1}y_{i,j}\prod_{b_j=0}\bar y_{i,j}$, such that the quadratic degree of $\Pi'$ is at most that of $\Pi$.
\begin{proof}
Note that the axiom $(\mathbf y_i\neq b_1\ldots b_{\log m})\cdot E_{b_1\ldots b_{\log m}}=0$ can be derived from the axiom $(\mathbf y_i\neq b_1\ldots b_{\log m})\equiv\prod_{b_j=1}y_{i,j}\prod_{b_j=0}\bar y_{i,j}$ by multiplying by the polynomial $E_{b_1\ldots b_{\log m}}$. Since this derivation involves only singular variables, the degree can never drop and therefore the quadratic degree of this derivation is at most that of the final polynomial. We construct $\Pi'$ as follows. We first derive the former axiom from the latter in $\Pi'$. Besides this derivation, $\Pi'$ involves the same steps as $\Pi$.
\end{proof}
\paragraph{Definition27.} An equation $E_{b_1\ldots b_{\log m}}$ is said to be banned when the previous lemma is applied repeatedly to eliminate all occurrences of it from the axioms.
\paragraph{Definition28.} A clean version of $F^{\mathsf{SEL}}_{n,k}$ is any subset of axioms of $F^{\mathsf{SEL}}_{n,k}$ along with axioms that ban some subset of equations of the form $E_{b_1\ldots b_{\log m}}$.
\paragraph{4.5.1.1 Cleanup($\rho$).}
We now describe how to perform the cleanup operations, which we collectively call $\mathsf{Cleanup}(\rho)$, that takes as input an ``unclean'' version of $F^{\mathsf{SEL}}_{n,k}$ derived by applying a restriction $\rho$ to a clean version, and outputs another clean version that is in some sense a subset of the input. Suppose that we are given a restriction $\rho$ that has been applied to a clean version of $F^{\mathsf{SEL}}_{n,k}$, with a variable $w^*\in\rho$ possibly set to $\star$, indicating that it must remain unset. To eliminate an axiom that has been affected by a $\mathbf x$ variable in $\rho$ not set to $\star$, we simply obtain the refutation that bans the corresponding equation $E_{b_1\ldots b_{\log m}}$ as described in the above lemma. Note that since we are eliminating the axiom without setting any variables in it, we can also do this in case our variable $w^*\in\mathbf x$. Suppose that $y_{ij}$ is a $\mathbf y$ variable in $\rho$ not set to $\star$. We first note that any axiom that contained $y_{ij}$ before the application of $\rho$ contains all the variables $y_{i1}\ldots y_{i\log m}$ corresponding to the $i$th pigeon $\mathbf y_i$. We first make sure that this $i$th pigeon does not contain our variable $w^*$ that must remain unset. If it doesn't, we proceed as follows. We set all the other variables in this pigeon to select some equation $E_{b_1\ldots b_{\log m}}$ that has not been banned. Such an equation exists provided that the number of banned equations so far is bounded, and the size of the restriction $\rho$ is also bounded (we formalize this in the lemma below). We then apply an additional restriction to the $\mathbf x$ variables that satisfies this equation $E_{b_1\ldots b_{\log m}}$ picked above. We then ban all the equations affected by this additional restriction, like we did above for the part of $\rho$ containing $\mathbf x$ variables. This eliminates the pigeon $\mathbf y_i$. We are left with the case where our variable $w^*$ belongs to some pigeon $\mathbf y_j$. We set all the variables in the pigeon $\mathbf y_j$ except for $w^*$, such that neither of the two equations $E_{b_1\ldots b_{\log m}}$ and $E_{b'_1\ldots b'_{\log m}}$ that would be selected if $w^*$ is set to zero or one are banned (again, these exist under the same conditions as above). We then proceed as before, i.e. apply an additional restriction to satisfy both these equations, and then ban any other equations that have been affected by this additional restriction. With this we have eliminated the axioms of pigeon $\mathbf y_j$ which select an equation, but we are still left with the axioms that prevent $\mathbf y_j$ from colliding with any other pigeon, which are now of the form $w^*\cdot(\mathbf y_{j'}\neq b_1\ldots b_{\log m})$ and $w^*\cdot(\mathbf y_{j'}\neq b'_1\ldots b'_{\log m})$ indicating that any pigeon $\mathbf y_{j'}$ distinct from $\mathbf y_j$ must not be mapped to the equations $E_{b_1\ldots b_{\log m}}$ and $E_{b'_1\ldots b'_{\log m}}$ if one of them is selected by setting $w^*$ to zero or one. To remove the latter axioms we do something similar to the process of banning an equation, where we simply replace these axioms by the axioms $(\mathbf y_{j'}\neq b_1\ldots b_{\log m})$ and $(\mathbf y_{j'}\neq b'_1\ldots b'_{\log m})$, effectively banning the equations $E_{b_1\ldots b_{\log m}}$ and $E_{b'_1\ldots b'_{\log m}}$ for the remaining pigeons.
\paragraph{4.5.1.2 Correctness of Cleanup($\rho$).}
We note that the above cleanup operations over $\mathbf y$ variables terminate successfully only when there are enough equations that have not been banned by prior calls to cleanup, and also the size of the restriction $\rho$ is bounded. We make this formal by the below lemma.
\paragraph{Lemma29 (Correctness of Cleanup($\rho$)).} Let $\rho$ be a restriction of size $\kappa$. If the number of banned equations (from previous calls to Cleanup) is $\ll m/2^{\kappa}$, then $\mathsf{Cleanup}(\rho)$ terminates correctly. Moreover, it bans at most $O(\kappa)$ additional equations and removes at most $O(\kappa)$ pigeons in its run.
\begin{proof}
In $\mathsf{Cleanup}(\rho)$, note that we can remove the axioms that contain $\mathbf x$ variables unconditionally. When we remove a pigeon $\mathbf y_i=y_{i1}\ldots y_{i\log m}$, we rely on having an equation it can be set to that is not already banned. Since the size of $\rho$ is bounded by $\kappa$, note that at most $\kappa$ variables from $y_{i1}\ldots y_{i\log m}$ can be set by $\rho$. Therefore there are at least $\log m-\kappa$ of them unset, corresponding to selecting $m/2^{\kappa}$ many equations. Since we assume that the number of banned equations is much less than this, we can always find one that is not banned to assign this pigeon to.
We now count the number of new equations banned and the number of pigeons removed by this call to $\mathsf{Cleanup}(\rho)$. Since each $\mathbf x$ variable appears in a constant number of equations, the number of equations we ban while processing it is a constant. When we process a $\mathbf y$ variable, we pick and satisfy an equation, and ban all other equations affected in the process. Since every equation also contains a constant number of variables, satisfying it affects only a constant number of other equations. Therefore, for every variable we process we ban only a constant number of equations, and thus the total number of equations banned is $O(\kappa)$. We remove only those pigeons with a variable in $\rho$, so this is also bounded by $O(\kappa)$.
\end{proof}
\subsubsection*{4.5.2 The Main Theorem}
We need first the following key lemma that shows how to apply the Split operation to reduce high quadratic degree terms.
\paragraph{Algorithm1: Algorithm for Lemma30.}
\textbf{Input:} A refutation $\Pi$, and a nonsingular variable $z$ with extension axiom $z-Q=0$ satisfying the pre-conditions of Lemma30.
\textbf{Output:} A refutation $\Pi'$ satisfying post-conditions of Lemma30.
\begin{enumerate}[label=\arabic*,leftmargin=*]
\item Let $\ell_0<\ell_1$ be such that $|H_d(\Pi,z,\ell_0,\ell_1)|\geq|H_d(\Pi,z)|/p^2$.
\item Apply Lemma22 with $l=\ell_1-\ell_0$ to obtain $\delta,w^*,a,b$ satisfying post-conditions of Lemma22.
\item $\Pi=\Pi|_{\delta}$ (and in particular $z-Q|_{\delta}=z-(b-a)w^*-a$).
\item $\mathsf{Cleanup}(\delta\cup\{w^*=\star\})$ (Cleanup axioms affected by $\delta$ and remove $w^*$ from all axioms other than $z-(b-a)w^*-a$ while keeping it alive.)
\item \textbf{if} $w^*$ contributes to $\geq\epsilon/4p^2$ fraction of pairs in $H_d(\Pi)$ \textbf{then}
\item $\Pi=\Pi|_{w^*=0}$
\item \textbf{end}
\item \textbf{else}
\item Apply the substitution $(z-a)/(b-a)$ for $w^*$ in $\Pi$.
\item Let $\Pi'=\Split_{z,l_0,l_1}(\Pi)$.
\item \textbf{end}
\end{enumerate}

\paragraph{Lemma30.} Let $F$ be a system of unsatisfiable polynomials and let $z$ be a nonsingular extension variable with the extension axiom $z-Q$. Let $\ell=\ell(\supp(Q))$ so that $z^{\ell}=c$ holds for some $c\in\mathbb F_p$. Let $\Pi$ be a refutation of $F\cup\{z-Q\}$ modulo $z^{\ell}=c$ such that for at least an $\epsilon$ fraction of pairs $(t_1,t_2)$ in $H_d(\Pi)$, $\Qdeg(t_1,t_2,z)=1$, for some $d\geq0$. Then there exists a refutation $\Pi'$ of $F$ such that $|H_d(\Pi')|\leq(1-\epsilon/4p^2)|H_d(\Pi)|$.
\begin{proof}
We will apply a procedure as described by Algorithm1 in order to modify the proof to satisfy the post-conditions of the Lemma. Here we give a detailed description of the algorithm, together with its correctness. Let $H_d(\Pi,z)$ be the set of all unordered pairs $(t_1,t_2)\in H_d(\Pi)$ that $z$ contributes to. That is, $H_d(\Pi,z)$ is the set of all unordered pairs $(t_1,t_2)\in H_d(\Pi)$ such that $\Qdeg(t_1,t_2,z)=1$. There are many different ways that $z$ can contribute to $H_d(\Pi,z)$: namely, for all $i,j$ such that $i<j<\ell$, let $H_d(\Pi,z,i,j)$ be the set of all unordered pairs $(t_1,t_2)\in H_d(\Pi,z)$, such that the degree of $z$ in $t_1$ is $i$ and the degree of $z$ in $t_2$ is $j$. Note that for any two pairs $(i,j)$ and $(i',j')$ such that $i\neq i'$ or $j\neq j'$, $H_d(\Pi,z,i,j)$ and $H_d(\Pi,z,i',j')$ are disjoint. Therefore, there exists a ``good'' pair $\ell_0<\ell_1<\ell$ such that removing $H_d(\Pi,z,\ell_1,\ell_0)$ from $H_d(\Pi,z)$ will remove at least a $1/p^2$ fraction of $H_d(\Pi,z)$ and therefore a $\epsilon/p^2$ fraction of pairs in $H_d(\Pi)$, since $|H_d(\Pi,z)|\geq\epsilon|H_d(\Pi)|$.
We want to apply the Split operation $\Split_{z,\ell_0,\ell_1}$ to remove all such pairs. But in order to do this we have to satisfy the preconditions of Lemmas23 and25: we need two values $a,b$ such that $a^{\ell_1-\ell_0}\neq b^{\ell_1-\ell_0}$ and all the axioms should be free of $z$ except for $(z-a)(z-b)=0$. The first step (Line2 of1) is to apply Lemma22 with $l=\ell_1-\ell_0$. This gives us $w^*\in\vars(Q)$, $a,b\in\supp(Q)$ and a partial restriction $\delta$ to $\vars(Q)-w^*$ such that $(z-Q)|_{\delta}=z-(b-a)w^*-a$, where $a^{\ell_1-\ell_0}\neq b^{\ell_1-\ell_0}\bmod p$. Next, we apply the restriction $\delta$ to $\Pi$ (Line3).
Now we have a simpler linear extension axiom for $z$ of the form $z-(b-a)w^*-a=0$. Next we would like to make the substitution $w^*=(z-a)/(b-a)$ in $\Pi$ in order to satisfy this extension axiom, towards the goal of eliminating $z$ from the axioms so that we have the preconditions of Lemma25 and therefore are able to apply $\Split_{z,\ell_1,\ell_0}$. However, if $w^*$ appears in any of the axioms in $F$, this would create additional occurrences of $z$ and we would not make any progress. Therefore, we have to make sure that none of the axioms of $F$ contain $w^*$. But we also cannot set $w^*$ to zero or one in an attempt to get rid of it, since this would set $z$ to either $a$ or $b$ through the above extension axiom, and Split requires that $z$ take on two distinct values. We thus have to get rid of all axioms mentioning $w^*$ either by setting other variables or by replacing these axioms with stronger versions, such that the former can be derived from the latter. This is what the subroutine Cleanup does, in addition to removing the axioms in $F$ that were affected by our earlier restriction $\delta$, so that we have a clean version of $F^{\mathsf{SEL}}_{n,k}$ as defined in the previous section.
We are now ready to make the substitution $w^*=(z-a)/(b-a)$. Under this substitution, the Boolean axiom $w^{*2}-w=0$ reduces to $(z-a)(z-b)=0$, and the original extension axiom for $z$ disappears (since under this substitution it becomes $0=0$.) Thus this substitution would satisfy all of the preconditions of Lemmas23,25. However, this substitution can create a new problem: it can cause a blow up in the size of $H_d(\Pi)$ since for every pair of terms $(t_1,t_2)$ such that one of them contains $w^*$, we could have up to four new terms after the substitution. In order to deal with this potential blow up we do a simple case analysis: If $w^*$ contributes to at least an $\epsilon/4p^2$ fraction of pairs $(t_1,t_2)$ in $H_d(\Pi)$, then we set $w^*=0$ (Lines4–5). This gives us the required reduction in the size of $H_d(\Pi)$ ($z$ is also set to a constant by setting $w^*=0$, but we don't care about that since we have obtained a reduction in high Quadratic degree terms without needing to use Split). Otherwise, the blowup caused by the substitution $w^*=(z-a)/(b-a)$ adds at most $3\epsilon/4p^2$ fraction of pairs to $H_d(\Pi)$, and thus if we remove all pairs in $H_d(\Pi,z,\ell_0,\ell_1)$ (after this blowup) then overall we will will have reduced the size of $H_d(\Pi)$ to $(1-\epsilon/4p^2)|H_d(\Pi)|$. So in this latter case, we apply the substitution mentioned above (Line8) which simultaneously removes $w^*$ from all axioms, and replaces the linear axiom for $z$ by $(z-a)(z-b)=0$. Now all preconditions for Lemma8 hold so we can apply $\Split_{z,\ell_0,\ell_1}$ (Line9) to get a valid refutation. It is left to argue that this indeed removes the set $H_d(\Pi,z,\ell_1,\ell_0)$. More precisely, we argue that high Quadratic degree pairs of terms in the refutation obtained after applying Split have a one to one mapping to the set $H_d(\Pi)-H_d(\Pi,z,\ell_1,\ell_0)$. Fix a line $P\in\Pi$. Since we are working modulo $z^{\ell}=c$, we can assume that $P=P_0+zP_1+\ldots z^{\ell-1}P_{\ell-1}$. Let $R=z^{\ell_0}R_0+z^{\ell_1}R_1$ be the unique polynomial equivalent to $P\bmod(z-a)(z-b)$. $\Split_{z,\ell_0,\ell_1}(\Pi)$ is the refutation with lines $R_1,R_0$ for all $P\in\Pi$. By the proof of Lemma23 $R_0,R_1$ have the form:
\[R_1=P_{\ell_1}+\sum_{i<\ell,i\neq\ell_0}c_{1i}P_i,\qquad R_0=P_{\ell_0}+\sum_{i<\ell,i\neq\ell_1}c_{0i}P_i\]
for some constants $c_{1i},c_{0i}\in\mathbb F_p$.
For a pair of terms $(t_i,t_j)$ in $R_1$ such that $t_i\in P_i$ and $t_j\in P_j$ and $\Qdeg(t_i,t_j)\geq d$, we map it to the pair $(t_iz^i,t_jz^j)\in P$, and similarly for $R_0$. Clearly this is a one-one mapping, and since $P_{\ell_0}$ does not occur in $R_1$ and $P_{\ell_1}$ does not occur in $R_0$, it is a mapping to $H_d(\Pi)-H_d(\Pi,z,\ell_1,\ell_0)$. Therefore we have that for the refutation $\Pi'=\Split_{z,\ell_0,\ell_1}(\Pi)$ whose lines are $\{R_1,R_0\}$,
\[|H_d(\Pi')|\leq|H_d(\Pi)-H_d(\Pi,z,\ell_1,\ell_0)|\leq(1-\epsilon/4p^2)|H_d(\Pi)|.\]
\end{proof}

\paragraph{Theorem31.} For $n$ sufficiently large, any $(M,\kappa)$-PC+Ext refutation of $F^{\mathsf{SEL}}_{n,k}$ has size
\[\exp\left(\frac{\Omega(n^2)}{10^{\kappa}(M+n\log n)}\right).\]

\begin{proof}
Let $\Pi$ be an alleged $(M,\kappa)$-PC+Ext refutation of $F^{\mathsf{SEL}}_{n,k}$ with logarithm of its size less than $\gamma n^2/(10^{\kappa}(M+n\log n))$, for a small enough constant $\gamma$. Given $\Pi$, Algorithm2 (defined below) will apply a sequence of restrictions and cleanup steps in order to produce a refutation $\Pi'$ of a clean version of $F^{\mathsf{SEL}}_{n,k}$ (see Definition29) with the property that the Quadratic degree of $\Pi'$ is at most $d=\nu n/\kappa$ for a small enough constant $\nu>0$. The algorithm contains a while loop which iteratively removes all pairs of terms of high Quadratic degree. From $\Pi'$, we will apply a further restriction to all of the remaining unset $\mathbf y$-variables (i.e. pigeons that select equations from $F_{n,k}$), to extract a refutation of a subset of $m'$ equations from $F_{n,k}$ of low degree, contradicting the degree lower bound given in Lemma35. Recall that $F_{n,k}$ is defined over variables $\mathbf x$ and we pick a subset of these equations by matching pigeons $\mathbf y_i$ to equations in $F_{n,k}$ through a complete bipartite graph.
The algorithm first initializes a few things. Set $d=\nu n/\kappa$ for a small enough constant $\nu>0$. Let $M'=M+n\log n$, which upper bounds the total number of variables occurring in the refutation. Let $S$ be the set of all variables that are Singular w.r.t. the current set of extension axioms. We initialize $S$ to be the set of all variables $\mathbf x\cup\mathbf y\cup\mathbf z$ since this is the largest possible set we will be dealing with; this will be updated at every iteration of the while loop, although we note that it can only reduce as we apply restrictions. Henceforth when we refer to Quadratic degree, we mean $\Qdeg_S$. Finally, we initialize $H$ to be the set of all pairs of terms in $\Pi$ with Quadratic degree greater than $d$.
In the while loop, we first update the set $S$ by checking which of the extension variables $z$ have zero in their support according to their current extension axioms, and deleting those that don't. For each extension variable $z$ that we delete from $S$, we reduce the refutation $\Pi$ by $z^{\ell(z)}=c$. Such an identity exists and is derivable by Lemma13, and does not increase the size of $H$ by Lemma20. Once we have updated $S$, we recompute the set of high Quadratic degree pairs $H$ with respect to the updated set $S$. This also does not increase the size of $H$, by Lemma16. We then pick a variable $w$ that contributes to the Quadratic degree of at least a $d/M'$ fraction of pairs in $H$ by averaging.
\paragraph{Algorithm2: Eliminating high Quadratic degree terms from the proof.}
\textbf{Input:} A refutation $\Pi$ of $F^{\mathsf{SEL}}_{n,k}$ with extension axioms $\mathsf{Ext}$.
\textbf{Output:} A refutation $\Pi'$ with Quadratic degree less than $d$.
\begin{enumerate}[label=\arabic*,leftmargin=*]
\item $d\leftarrow\nu n/\kappa$, where $\nu$ is a sufficiently small constant.
\item $M'\leftarrow M+n\log(n)$. ($M'$ upper bounds $|\mathbf x\cup\mathbf y\cup\mathbf z|$, the total number of variables.)
\item $S\leftarrow\mathbf x\cup\mathbf y\cup\mathbf z$ (the current set of singular variables: all Boolean variables are singular by default and we initialize all extension variables to also be singular. This could possibly reduce in each iteration.)
\item $H\leftarrow\{(t_1,t_2)\mid t_1,t_2\in\Pi\text{ and }\Qdeg_S(t_1,t_2)\geq d\}$ (the set of all pairs of terms of large Quadratic degree according to $S$).
\item \textbf{while} $H$ is non empty \textbf{do}
\item \textbf{for every} extension axiom $z-Q\in\mathsf{Ext}$ \textbf{do}
\item \textbf{if} $0\notin\supp(Q)$ \textbf{then}
\item $S\leftarrow S-\{z\}$
\item Compute $c$ such that $z^{\ell(z)}=c$ and reduce $\Pi$ by the latter identity.
\item \textbf{end}
\item \textbf{end}
\item $H\leftarrow\{(t_1,t_2)\mid t_1,t_2\in\Pi\text{ and }\Qdeg_S(t_1,t_2)\geq d\}$ (update $H$ to reflect changes due to the above for loop).
\item Pick a variable $w$ that, by an averaging argument, occurs in at least an $\epsilon$ fraction of terms in $H$, where we choose $\epsilon=d/M'$.
\item \textbf{if} $w\in S$ \textbf{then}
\item Let $\sigma$ be a restriction on $\mathbf x\cup\mathbf y$ such that $w|_{\sigma}=0$.
\item $\Pi\leftarrow\Pi|_{\sigma}$
\item $\mathsf{Cleanup}(\sigma)$
\item \textbf{end}
\item \textbf{else}
\item Apply Algorithm1, which by Lemma30 satisfies the post-conditions of Lemma30.
\item \textbf{end}
\item \textbf{end}
\end{enumerate}
There are two cases depending on whether $w\in S$ or not. In the first case (lines14–18), $w$ is Singular so we apply the restriction $\sigma$ such that $w|_{\sigma}=0$ and call $\mathsf{Cleanup}(\sigma)$ to restore to a clean version of our tautology. This eliminates the contribution to high Quadratic degree from terms containing $w$, and hence obtains a $(1-d/M')$-factor reduction in the size of $H$. In the second case (lines19–34), $w$ is Nonsingular so we apply Algorithm1, which uses the Split operation non-trivially to reduce the size of $H$. Lemma30 proves correctness of the algorithm, and thus upon termination of one call to Algorithm1, we have obtained a $(1-d/(4p^2M'))$-factor reduction in the number of high Quadratic degree terms.
Repeating the above for
\[\frac{-\log|H|}{\log(1-d/4p^2M')}\approx\frac{4p^2M'\log|H|}{d}\leq O(\gamma)\kappa n/10^{\kappa}\]
iterations, we eliminate all terms in $H$ from the proof and thus obtain a refutation of Quadratic degree less than $d$. Since we call Cleanup once per iteration, and in each call it bans at most $O(\kappa)$ many equations and removes at most $O(\kappa)$ many pigeons (by Lemma29), we have banned at most $O(\gamma)\kappa^2n/10^{\kappa}$ equations and removed at most those many pigeons in total. Therefore, we always satisfy the invariant that the number of banned equations is much less than $m/2^{\kappa}$ (where $m=10n$), satisfying the required conditions for correctness of Cleanup from Lemma29.
Let $\Pi'$ denote the modified proof upon termination of Algorithm2. Note that out of the $m'=(1-\epsilon)m$ pigeons, there are at least a $1-O(\gamma)$ fraction of pigeons still alive (i.e. not removed by Cleanup) and a $1-O(\gamma)$ fraction of the $m$ equations not banned. We now substitute for the remaining pigeons $\mathbf y$ so that we select a subset of at least $(1-2\epsilon)m$ unsatisfiable equations from $F_{n,k}$ that are not banned and obtain a refutation of them of Quadratic degree at most $d$ (assuming $\gamma$ is small enough). By Lemma21, we can obtain a refutation of these equations of degree at most $3pd$. Now, for all surviving extension variables we substitute them with their definitions in terms of the variables $\mathbf x$. Note that since each extension variable is a degree $\kappa$ polynomial this raises the degree to at most $3\kappa pd$. Since $d=\nu n/\kappa$, for sufficiently small $\nu$ we end up with a refutation of $(1-2\epsilon)$ equations from $F_{n,k}$ of degree less than $c_2n$, contradicting Lemma35.
\end{proof}

\section*{A Appendix}
We will prove Theorem9, which we state again here for convenience.
\paragraph{Theorem32 (Theorem9).} Let $m=10n$. Then there exists constants $k>0$, $0<\epsilon<1$ such that for sufficiently large $n$, there exists $k$-CSP formulas $\{F_{n,k}\}$ with $m$ $k$-local constraints such that for $m'=(1-\epsilon)m$, every subset of $m'$ constraints is unsatisfiable and requires linear degree PC refutations.
First we'll show that a random regular bipartite graph has good boundary expansion. This has been used implicitly in other works ([9],[5]), but for completeness we state and prove it here. Let $G=(L,R,E)$ be a bipartite graph, and let $A\subseteq R$. The boundary for $A$, $\partial(A)$, is the set of vertices $x$ in $L$ so that $|N(x)\cap A|=1$, i.e., vertices with a unique neighbor in $A$. A bipartite graph is $(d,k)$ regular if every vertex in $L$ has degree $d$ and every vertex in $R$ has degree $k$. In this case, for $n=|L|$, $m=|R|$, we have $dn=km$.
\paragraph{Theorem33.} Let $d,k,n,m$ be positive integers with $dn=km$, $k\geq12$. Then with high probability for a random $(d,k)$ regular bipartite graph with $|L|=n$, $|R|=m$, for all $A\subset R$, $|A|<n/(e^6k^2)$, we have $\partial(A)\geq k|A|/2$.
\begin{proof}
Let $N(A)$ be all the neighbors of $A$. Since the total degrees of vertices in $A$ is $k|A|$, and each element of $N(A)-\partial(A)$ is contingent on two such edges, $k|A|\geq2(|N(A)|-|\partial(A)|)+|\partial(A)$, or $\partial(A)\geq2|N(A)|-k|A|$. We will show that with high probability for all such $A$, $|N(A)|>3k|A|/4$, and hence $\partial(A)\geq k|A|/2$.
If not, there are sets $A\subset R$ and $B\subset L$ so that $N(A)\subseteq B$ and $|B|=3k|A|/4$. We will bound the probability that this is true for fixed sets $A,B$ and then take a union bound. We can view picking a random $(d,k)$ bipartite graph as picking a random matching between $d$ half-edges adjacent to each $x\in L$ and $k$ such half-edges adjacent to each $y\in R$; if a half edge for $x$ is matched to a half-edge for $y$, it forms an edge between $x$ and $y$.
We can form this matching by going through the half edges for nodes in $R$ and for each randomly selecting an unmatched half-edge for some node in $L$. We start with the edges for $A$ in an arbitrary order. If we condition on all previous neighbors for $A$ being in $B$, the number of half-edges left still available for $B$ is less than $d|B|$, whereas the number for $\bar B$ stays at exactly $d(n-|B|)$. Thus, the conditional probability that the next edge formed is also in $B$ is at most $|B|/n$, and we do this for each of $k|A|$ edges, meaning the probability that all neighbors are in $B$ is at most $(|B|/n)^{k|A|}$.
Now, for a fixed $|A|$ and setting $|B|=3k|A|/4$, we take the union bound over all subsets $A$ and $B$. This gives a total probability of failure for some set $A$ of size $a$ as:
\begin{align*}
&\binom ma\binom n{3ka/4}(3ka/4n)^{ka}\\
&\leq(em/a)^a(4en/3ka)^{3ka/4}(3ka/4n)^{ka}\\
&\leq(em/a)^a(e^3ka/n)^{ka/4}\\
&=(ekn/da)^a(e^3ka/n)^{ka/4}\\
&=\left(\frac{e^{3k/4+1}a^{k/4-1}k^{k/4+1}}{dn^{k/4-1}}\right)^a.
\end{align*}
Since we are assuming $a<n/(e^6k^2)$, the base in the above expression is at most
\[\frac{e^{3k/4+1}(n/e^6k^2)^{k/4-1}k^{k/4+1}}{dn^{k/4-1}}=e^{7-3k/4}k^{3-k/4}/d\]
which for $k\geq12$ is bounded below $e^{-2}$, meaning the probability of such a bad set existing is exponentially small in $a$, and the probability of such a bad set existing for any $a$ is less than $1/2$.
\end{proof}
\paragraph{Definition34.} For a Boolean vector $X=\{x_1,\ldots,x_n\}$, we define $\mathcal L_{n,m,k_1,k}(X)$ to be the distribution over $k$-CSP formulas over $n$ variables $X=\{x_1,\ldots,x_n\}$ obtained by selecting $m$ parity equations, where each parity is represented by a node on the right of a randomly chosen bipartite graph $G(L,R,E)$, with $|L|=n$, $|R|=m$, and with left degree bounded by $k_1$ and right degree bounded by $k$.
\paragraph{Lemma35.} Let $F_{n,k}$ be a tautology given by the system of parity equations $AX=b$ over variables $X=\{x_1,\ldots,x_n\}$ drawn at random from $\mathcal L_{n,m,k_1,k}$ where $m=10n$, for large enough constants $k_1,k>0$, and $b$ is chosen randomly. Then the following hold with high probability for a small enough $\epsilon>0$:
\begin{enumerate}[label=\alph*)]
\item Any subset of a $(1-\epsilon)$-fraction of the equations in $F_{n,k}$ is unsatisfiable.
\item Any subset of a $(1-\epsilon)$-fraction of the equations in $F_{n,k}$ requires PC degree $c_2n$ to refute, for some $c_2>0$.
\end{enumerate}
\begin{proof}
\begin{enumerate}[label=\alph*)]
\item The probability that a set of $(1-\epsilon)10n$ random parities (i.e. for a random choice of $b$) is satisfiable is at most $2^{-9n}$ for a small enough $\epsilon$. The probability that any such subset of $F_{n,k}$ is satisfiable is therefore at most $2^{(-n(9-10H(\epsilon)))}$, which is exponentially small for a small enough $\epsilon$ (where $H(\epsilon)$ is the binary entropy function).
\item This follows directly from [1], Theorem3.8 and Theorem4.4, since by Theorem33 the bipartite graph underlying the system of parity equations $A$ has good boundary expansion with high probability.
\end{enumerate}
\end{proof}

\begin{thebibliography}{99}
\bibitem[1]{ref1} Michael Alekhnovich and Alexander A. Razborov. Lower bounds for polynomial calculus: Non-binomial case. In 42nd Annual Symposium on Foundations of Computer Science, FOCS 2001, 14-17 October 2001, Las Vegas, Nevada, USA, pages 190–199. IEEE Computer Society, 2001. doi:10.1109/SFCS.2001.959893.
\bibitem[2]{ref2} Yaroslav Alekseev. A lower bound for polynomial calculus with extension rule. In Valentine Kabanets, editor, 36th Computational Complexity Conference, CCC 2021, July 20-23, 2021, Toronto, Ontario, Canada (Virtual Conference), volume 200 of LIPIcs, pages 21:1–21:18. Schloss Dagstuhl – Leibniz-Zentrum für Informatik, 2021. doi:10.4230/LIPIcs.CCC.2021.21.
\bibitem[3]{ref3} Robert Andrews and Michael A. Forbes. Ideals, determinants, and straightening: Proving and using lower bounds for polynomial ideals. CoRR, abs/2112.00792, 2021. arXiv:2112.00792.
\bibitem[4]{ref4} Paul Beame, Russell Impagliazzo, Jan Krajíček, Toniann Pitassi, and Pavel Pudlák. Lower bounds on hilbert’s nullstellensatz and propositional proofs. Proceedings of the London Mathematical Society, 3(1):1–26, 1996.
\bibitem[5]{ref5} Paul Beame, Richard Karp, Toniann Pitassi, and Michael Saks. The efficiency of resolution and davis–putnam procedures. SIAM Journal on Computing, 31(4):1048–1075, 2002.
\bibitem[6]{ref6} Eli Ben-Sasson and Russell Impagliazzo. Random cnfs are hard for the polynomial calculus. In 40th Annual Symposium on Foundations of Computer Science (Cat. No. 99CB37039), pages 415–421. IEEE, 1999.
\bibitem[7]{ref7} Eli Ben-Sasson and Avi Wigderson. Short proofs are narrow – Resolution made simple. In Proceedings of the thirty-first annual ACM symposium on Theory of computing, pages 517–526, 1999.
\bibitem[8]{ref8} Sam Buss, Dima Grigoriev, Russell Impagliazzo, and Toniann Pitassi. Linear gaps between degrees for the polynomial calculus modulo distinct primes. Journal of Computer and System Sciences, 62(2):267–289, 2001.
\bibitem[9]{ref9} Vašek Chvátal and Endre Szemerédi. Many hard examples for resolution. Journal of the ACM (JACM), 35(4):759–768, 1988.
\bibitem[10]{ref10} Matthew Clegg, Jeffery Edmonds, and Russell Impagliazzo. Using the groebner basis algorithm to find proofs of unsatisfiability. In Proceedings of the twenty-eighth annual ACM symposium on Theory of computing, pages 174–183, 1996.
\bibitem[11]{ref11} Stephen A Cook and Robert A Reckhow. The relative efficiency of propositional proof systems. The journal of symbolic logic, 44(1):36–50, 1979. 7:23
\bibitem[12]{ref12} Stefan S. Dantchev and Søren Riis. On relativisation and complexity gap. In Matthias Baaz and Johann A. Makowsky, editors, Computer Science Logic, 17th International Workshop, CSL 2003, 12th Annual Conference of the EACSL, and 8th Kurt Gödel Colloquium, KGC 2003, Vienna, Austria, August 25-30, 2003, Proceedings, volume 2803 of Lecture Notes in Computer Science, pages 142–154. Springer, 2003. doi:10.1007/978-3-540-45220-1\_14.
\bibitem[13]{ref13} Michael A. Forbes, Amir Shpilka, Iddo Tzameret, and Avi Wigderson. Proof complexity lower bounds from algebraic circuit complexity. Theory Comput., 17:1–88, 2021. URL: https://theoryofcomputing.org/articles/v017a010/.
\bibitem[14]{ref14} Nashlen Govindasamy, Tuomas Hakoniemi, and Iddo Tzameret. Simple hard instances for low-depth algebraic proofs. In 63rd IEEE Annual Symposium on Foundations of Computer Science, FOCS 2022, Denver, CO, USA, October 31 – November 3, 2022, pages 188–199. IEEE, 2022.
\bibitem[15]{ref15} Armin Haken. The intractability of resolution. Theoretical computer science, 39:297–308, 1985.
\bibitem[16]{ref16} Russell Impagliazzo, Sasank Mouli, and Toniann Pitassi. The surprising power of constant depth algebraic proofs. In Proceedings of the 35th Annual ACM/IEEE Symposium on Logic in Computer Science, pages 591–603, 2020.
\bibitem[17]{ref17} Russell Impagliazzo, Pavel Pudlák, and Jirí Sgall. Lower bounds for the polynomial calculus and the gröbner basis algorithm. Comput. Complex., 8(2):127–144, 1999. doi:10.1007/ s000370050024.
\bibitem[18]{ref18} Alexander A Razborov. Lower bounds on the size of bounded depth circuits over a complete basis with logical addition. Math. notes of the Academy of Sciences of the USSR, 41(4):333–338, 1987.
\bibitem[19]{ref19} Roman Smolensky. Algebraic methods in the theory of lower bounds for boolean circuit complexity. In Proceedings of the nineteenth annual ACM symposium on Theory of computing, pages 77–82, 1987.
\bibitem[20]{ref20} Dmitry Sokolov. (semi) algebraic proofs over \{±1\} variables. In Proceedings of the 52nd Annual ACM SIGACT Symposium on Theory of Computing, pages 78–90, 2020.
\end{thebibliography}
\end{document}
